% !PRINTED 325
%<<CONT
und den Werth $X:Y$ durch einen variablen Punkt dieses Kreises
darstellen, und die betreffenden Punkte mit $\alpha, x, X$ bezeichnen, so
geht $\Phi$ durch Null, wenn $X$ von $x$ aus nach vorwärts oder nach
rückwärts bis zu den nächst gelegenen Punkten, die wir mit $\alpha_1$ und
$\alpha_2$ bezeichnen wollen, geht. Sind also $X_1$, $X_2$ die Wurzeln der
Gleichung (9), so haben wir folgende Anordnung auf dem Kreise

\weq{10}{\alpha_1,\ X_1,\ x,\ X_2,\ \alpha_2,}

\gesperrt{und $(X_1, X_2)$ stellt ein Intervall dar, in dem zwar der beliebig
gewählte Punkt $x$, aber kein Wurzelpunkt der Gleichung $f = 0$
liegt. $X_1$ und $X_2$ sind also den beiden nächst gelegenen Wurzeln
$\alpha_1, \alpha_2$ näher als der Ausgangswerth $x$.}

Dies gilt, welche Lage auch der zweite willkürliche Punkt $\xi$
haben mag. Es werden davon zwar die Punkte $X_1, X_2$ abhängig
sein; immer aber werden sie in dem Intervall $(\alpha_1, x)$ und $(x, \alpha_2)$
liegen.

Es entsteht nun die Frage: wie ist der Punkt $\xi$ zu wählen, damit
$X_1$ möglichst nahe an $\alpha_1$ oder $X_2$ möglichst nahe an $\alpha_2$ liege,
oder vielmehr, da es auf die Kenntniss von $\xi$ selbst nicht ankommt,
welches ist der möglichst nahe bei $\alpha_1$ gelegene Punkt $X_1$ und der
möglichst nahe bei $\alpha_2$ gelegene Punkt $X_2$?

Denken wir uns den Punkt $X$ gegeben, so ist (9) eine
quadratische Gleichung für den Punkt $\xi$; zu jeder Lage von $X$
gehören also zwei Lagen von $\xi$. Aber nur solche Werthe von
$X_1$ oder $X_2$ sind in (10) zulässig, für die diese quadratische
Gleichung reelle Wurzeln ergiebt. Also nur solche Werthe von
$X_1$ und $X_2$ können vorkommen, bei denen die Discriminante von
(9) in Bezug auf $\xi$ positiv ist, und wir erhalten die Grenzlagen
von $X_1$ und $X_2$, wenn wir die Discriminante von (9) gleich
Null setzen. Die Coëfficienten von $\xi^2$, $2\xi\eta$, $\eta^2$ in (9) sind aber
\[
\frac{f'(x)^2 - y^2 H}{nf^2}(Xy - Yx)^2 - Y^2
\]
\[
\frac{f'(x)f'(y) + xyH}{nf^2}(Xy - Yx)^2 + XY
\]
\[
\frac{f'(y)^2 - x^2 H}{nf^2}(Xy - Yx)^2 - X^2,
\]
und die Discriminante erhält man, wenn man von dem Quadrat
%CONT>>
